\documentclass[11pt,a4paper]{article}

\usepackage[utf8]{inputenc}
\usepackage[T1]{fontenc}
\usepackage[french]{babel}
\usepackage{lmodern}
\usepackage{geometry}
\usepackage{graphicx}
\usepackage{booktabs}
\usepackage{array}
\usepackage{enumitem}
\usepackage{fancyhdr}
\usepackage{lastpage}
\usepackage{titlesec}
\usepackage{hyperref}
\usepackage{xcolor}

\geometry{top=2.5cm,bottom=2.5cm,left=2.7cm,right=2.7cm}
\definecolor{ensablue}{RGB}{0,70,140}
\hypersetup{
    colorlinks=true,
    linkcolor=ensablue,
    urlcolor=ensablue,
    pdftitle={Rapport d'avancement — Semaine 1},
    pdfauthor={Groupe 3}
}

\pagestyle{fancy}
\fancyhf{}
\fancyhead[L]{\small Projet 3 — Théorie des jeux et SMA}
\fancyhead[R]{\small Rapport S1}
\fancyfoot[C]{\small Groupe 3 \quad | \quad Page \thepage/\pageref{LastPage}}
\setlength{\headheight}{14pt}

\titleformat{\section}
    {\Large\bfseries\color{ensablue}}
    {\thesection.}{0.6em}{}
\titleformat{\subsection}
    {\large\bfseries\color{ensablue}}
    {\thesubsection.}{0.6em}{}

\setlist[itemize]{leftmargin=1.2cm,itemsep=0.25em}
\setlist[enumerate]{leftmargin=1.2cm,itemsep=0.35em}

\begin{document}

% ============================================
% PAGE DE GARDE
% ============================================
\begin{titlepage}
    \centering
    \vspace*{0.3cm}

    % Placez ensalogo.png dans le même dossier que ce fichier .tex.
    \IfFileExists{ensalogo.png}{%
        \includegraphics[width=12cm]{ensalogo.png}%
    }{%
        \fbox{\parbox[c][3cm][c]{12cm}{\centering
        \textbf{École Nationale des Sciences Appliquées de Safi}\\
        Logo ENSA Safi à ajouter : \texttt{ensalogo.png}}}%
    }

    \vspace{0.5cm}

    \normalsize
    \textbf{Département IRT}\\
    \textbf{Génie Informatique et Intelligence Artificielle}\\
    \textbf{Option : IDIA}

    \vspace{1.2cm}

    \large
    \textbf{Rapport d'avancement — Semaine 1}

    \vspace{0.8cm}

    \Large
    \textbf{\color{ensablue}
    PROJET 3 : INFORMATION IMPARFAITE,\\
    ÉQUILIBRE BAYÉSIEN PARFAIT ET ENCHÈRES}

    \vspace{1.2cm}

    \normalsize
    \textbf{Réalisé par :}\\[0.2cm]
    ELATIFI Anas\\
    ELMAHTAL Mohamed Yassine\\
    ERRAFAI Achraf\\
    ERRAOUI Wiame

    \vspace{0.8cm}

    \textbf{Encadré par :}\\[0.2cm]
    M. Jraifi

    \vfill

    \normalsize
    \textbf{Module : Théorie des jeux et systèmes multi-agents}\\[0.2cm]
    \textbf{Année Universitaire : 2025 / 2026}
\end{titlepage}

\pagenumbering{roman}
\section*{Résumé exécutif}
\addcontentsline{toc}{section}{Résumé exécutif}

Ce rapport présente l'avancement du Groupe 3 au terme de la première semaine
du Projet 3 consacré à l'information imparfaite, à l'équilibre bayésien parfait
(EBP) et aux enchères. La période couverte est strictement limitée aux
journées des 11 et 12 novembre 2026. L'objectif de cette phase n'était pas
encore de produire un solveur complet, mais de construire une base de travail
cohérente : organiser le dépôt, préparer l'environnement logiciel, extraire
les notions théoriques essentielles et définir une trajectoire de développement
reproductible.

Le résultat principal de la Semaine 1 est un dépôt structuré, documenté et
adapté à un matériel limité. Le modèle local \texttt{llama3.2:1b} a été retenu
pour réduire les besoins en ressources. Les interfaces Python, les notebooks,
les tests hors ligne et la fiche de lecture sont présents sous forme de
préparatifs. Les fonctionnalités de calcul et de simulation restent
explicitement réservées aux semaines suivantes.

\clearpage
\tableofcontents
\clearpage
\pagenumbering{arabic}

\section{Introduction et contexte}

Les systèmes multi-agents prennent souvent des décisions alors qu'une partie
de l'information pertinente est privée, incomplète ou observée indirectement.
Dans ce contexte, la théorie des jeux fournit un cadre formel pour représenter
les interactions entre agents, leurs types, leurs croyances et leurs
stratégies. Le Projet 3 s'intéresse plus particulièrement aux jeux de
signalisation, à la notion d'équilibre bayésien parfait et aux mécanismes
d'enchères.

Le sujet combine ainsi une dimension théorique et une dimension
algorithmique. La première consiste à comprendre comment les croyances sont
mises à jour après l'observation d'un message et comment la rationalité
séquentielle est vérifiée. La seconde consiste à transformer ces notions en
fonctions Python testables, puis à comparer les résultats numériques aux
références théoriques.

La Semaine 1 joue le rôle de phase de cadrage. Elle permet d'éviter une
implémentation prématurée et de fixer les conventions communes du groupe :
organisation des fichiers, choix des dépendances, modèle local, format de
documentation et méthode de validation.

\section{Objectifs de la Semaine 1}

Les objectifs retenus pour les 11 et 12 novembre 2026 étaient les suivants :

\begin{enumerate}
    \item créer un dépôt et une arborescence conformes à l'Annexe A ;
    \item préparer un environnement Python léger et reproductible ;
    \item configurer Ollama avec un modèle compatible avec les ressources
    matérielles disponibles ;
    \item rédiger une fiche de lecture Phase 0 contenant les définitions,
    les formules et l'approche de calcul ;
    \item analyser les interfaces nécessaires au futur prototype ;
    \item formaliser le suivi de l'avancement dans un diagramme Gantt et un
    rapport hebdomadaire.
\end{enumerate}

\section{Travail réalisé depuis la semaine précédente}

Les travaux réalisés pendant la première semaine, les 11 et 12 novembre 2026,
sont les suivants :

\begin{itemize}
    \item création du dépôt GitHub et de la structure locale conformément à
    l'Annexe A ;
    \item installation et configuration d'Ollama avec le modèle léger
    \texttt{llama3.2:1b}, adapté à un ordinateur peu puissant ;
    \item rédaction de la fiche de lecture Phase 0 en LaTeX dans
    \texttt{docs/fiche\_lecture.tex} ;
    \item revue initiale du code de cours et préparation des interfaces Python ;
    \item création du diagramme Gantt de la Semaine 1 et du présent rapport.
\end{itemize}

\subsection{Organisation du dépôt}

La structure mise en place sépare clairement la documentation, les notebooks,
le code source et les tests. Le dossier \texttt{docs/} contient notamment la
fiche de lecture, le Gantt et les rapports d'avancement. Le dossier
\texttt{src/} contient les interfaces préparatoires \texttt{jeux.py},
\texttt{experiences.py} et \texttt{agent\_tj\_sma.py}. Enfin, les notebooks
servent à isoler les expérimentations pédagogiques et les tests permettent
de vérifier que la base du projet peut être exécutée hors ligne.

\subsection{Préparation de l'environnement}

Le fichier \texttt{requirements.txt} regroupe les bibliothèques nécessaires
au projet : calcul numérique avec NumPy et SciPy, manipulation de données avec
Pandas, visualisation avec Matplotlib, tests avec Pytest et intégration locale
avec LangChain et LangChain--Ollama. Le choix de ne pas ajouter de dépendances
non indispensables contribue à garder l'installation simple.

Le modèle Ollama est configuré par la variable d'environnement
\texttt{OLLAMA\_MODEL}, avec \texttt{llama3.2:1b} comme valeur par défaut.
Cette configuration évite de coder en dur un modèle différent et facilite la
reproduction de l'environnement par chaque membre du groupe.

\subsection{Préparation théorique}

La fiche de lecture présente les éléments nécessaires à l'implémentation :

\begin{itemize}
    \item la mise à jour bayésienne des croyances après l'observation d'un
    message ;
    \item la meilleure réponse du récepteur compte tenu de ses croyances ;
    \item la vérification des déviations profitables de chaque type de
    l'émetteur ;
    \item la distinction entre stratégies séparantes et regroupantes ;
    \item les repères théoriques des enchères au premier et au second prix.
\end{itemize}

Pour un message $m$ observé avec une probabilité positive, la croyance
conditionnelle s'écrit :
\[
\mu(t\mid m)=
\frac{\pi(t)\Pr(m\mid t)}
{\sum_{t'}\pi(t')\Pr(m\mid t')}.
\]
Le récepteur choisit ensuite une action maximisant son utilité espérée :
\[
a^*(m)\in\arg\max_{a\in A}
\sum_t \mu(t\mid m)u_R(t,m,a).
\]
Ces formules serviront de base au solveur prévu en Semaine 2.

\subsection{Préparation de l'agent}

Le fichier \texttt{agent\_tj\_sma.py} constitue un template d'intégration
locale. Il expose le nom du modèle configuré et réserve l'assemblage des outils
LangChain à une étape ultérieure. Cette séparation est volontaire : elle
permet de finaliser d'abord les fonctions de théorie des jeux et de ne pas
lier les tests fondamentaux à la disponibilité du serveur Ollama.

\section{Écarts entre le prévisionnel et le réalisé}

La réalisation est alignée à 100\,\% sur les objectifs de la Semaine 1.
Les tâches indiquées comme réalisées dans le diagramme Gantt correspondent
uniquement aux journées du 11/11/2026 et du 12/11/2026. Les activités des
semaines suivantes restent hors du périmètre réalisé de ce rapport.

\subsection{Lecture du diagramme}

Dans le Gantt, seules les cinq tâches de la section Semaine 1 portent
l'attribut \texttt{done}. Les dates affichées sont le 11/11/2026 et le
12/11/2026. Aucune tâche de prototype, de simulation, de test avancé ou de
rapport final n'est présentée comme réalisée. Cette distinction permet de
séparer clairement l'état réel du dépôt du planning prévisionnel.

\section{Difficultés et blocages rencontrés}

La principale contrainte concerne les ressources matérielles limitées. Le choix
du modèle local \texttt{llama3.2:1b} réduit la consommation mémoire et permet
de préparer l'agent dans un environnement adapté à un ordinateur peu puissant.
Les tests hors ligne restent indépendants du service Ollama, ce qui facilite
la validation locale de la structure du projet.

\subsection{Contraintes matérielles}

L'utilisation d'un modèle local implique de tenir compte de la mémoire
disponible, du temps de réponse et de la capacité du poste à exécuter
simultanément l'environnement Python et Ollama. Le choix de
\texttt{llama3.2:1b} constitue donc un compromis entre légèreté et utilité
pour les tâches d'assistance. Les expérimentations futures devront limiter
la taille des entrées, contrôler le nombre d'itérations et conserver des
graines aléatoires afin de garantir des résultats comparables.

\subsection{Mesures de réduction des risques}

Plusieurs mesures ont été retenues dès S1 :

\begin{itemize}
    \item conserver des tests hors ligne ne nécessitant pas Ollama ;
    \item utiliser des stubs explicites plutôt que des implémentations
    incomplètes présentant un faux résultat ;
    \item isoler la configuration du modèle dans une variable d'environnement ;
    \item documenter les commandes d'installation et d'exécution dans le
    fichier \texttt{README.md}.
\end{itemize}

\section{Livrables produits en S1}

\begin{table}[htbp]
\centering
\caption{Livrables effectivement produits pendant la Semaine 1}
\begin{tabular}{p{0.35\textwidth}p{0.52\textwidth}}
\toprule
\textbf{Élément} & \textbf{Contenu} \\
\midrule
\texttt{requirements.txt} & Dépendances Python minimales du projet. \\
\texttt{README.md} & Installation, configuration Ollama et exécution des tests. \\
\texttt{docs/fiche\_lecture.tex} & Notions EBP, enchères, formules et approche de calcul. \\
\texttt{docs/gantt\_previsionnel.mmd} & Source Mermaid du suivi S1. \\
\texttt{docs/gantt\_previsionnel.png} & Rendu graphique du Gantt S1. \\
\texttt{src/} & Interfaces préparatoires sans calcul déclaré comme finalisé. \\
\texttt{tests/} & Vérification minimale de l'exécution hors ligne. \\
\bottomrule
\end{tabular}
\end{table}

\section{Vérification de la base technique}

La base S1 a été vérifiée sur trois aspects. Premièrement, les notebooks sont
des fichiers JSON valides et peuvent être ouverts par un environnement
Jupyter. Deuxièmement, les modules Python préparatoires sont syntaxiquement
compilables. Troisièmement, les tests hors ligne s'exécutent sans démarrer
Ollama. Ces contrôles ne constituent pas encore une validation scientifique
du solveur : ils confirment uniquement que la base de développement est
cohérente et prête à être étendue.

\section{Critères de réussite de la Semaine 1}

La Semaine 1 est considérée comme réussie selon les critères suivants :

\begin{itemize}
    \item le périmètre S1 est documenté et distinct des travaux ultérieurs ;
    \item le dépôt possède une organisation compréhensible par les quatre
    membres ;
    \item l'installation des dépendances et la configuration du modèle sont
    décrites ;
    \item les principales équations et notions du projet sont accessibles
    dans la fiche de lecture ;
    \item le suivi des tâches réalisées les 11 et 12 novembre est visible
    dans le Gantt ;
    \item les fichiers préparatoires ne produisent pas de résultats fictifs
    avant l'implémentation du solveur.
\end{itemize}

\section{Plan d'action pour la Semaine 2}

La Semaine 2 sera consacrée au prototype Python et au solveur :

\begin{enumerate}
    \item formaliser le jeu de signalisation ;
    \item implémenter le prototype des utilités et des croyances ;
    \item développer le solveur \texttt{equilibres\_ebp} ;
    \item préparer les premiers tests unitaires.
\end{enumerate}

\section{État du dépôt Git}

Le dépôt contient les fichiers de cadrage S1, les dépendances, les stubs
Python, les notebooks, la fiche de lecture et la documentation. Les commits
suivent la convention \texttt{[Prénom] verbe + objet}, par exemple :
\texttt{[Anas] rédiger la fiche de lecture Phase 0}.

\section{Répartition des rôles}

\begin{center}
\begin{tabular}{>{\raggedright\arraybackslash}p{0.38\textwidth}
                >{\raggedright\arraybackslash}p{0.48\textwidth}}
\toprule
\textbf{Membre} & \textbf{Rôle} \\
\midrule
ELATIFI Anas & Chef de projet / Théorie \\
ELMAHTAL Mohamed Yassine & Quant / Python \\
ERRAFAI Achraf & Développeur agent \\
ERRAOUI Wiame & Testeur / Expérimentateur \\
\bottomrule
\end{tabular}
\end{center}

\section{Diagramme Gantt de la Semaine 1}

Le diagramme ci-dessous présente uniquement les tâches réalisées les
11/11/2026 et 12/11/2026.

\begin{figure}[htbp]
    \centering
    \includegraphics[width=\textwidth]{gantt_previsionnel.png}
    \caption{Planification et avancement réalisés pendant la Semaine 1.}
    \label{fig:gantt-s1}
\end{figure}

\section{Conclusion}

La première semaine a permis de passer d'un sujet général à une base de projet
organisée et documentée. Le groupe dispose désormais d'un dépôt structuré,
d'une fiche théorique, d'un environnement léger et d'un planning cohérent avec
les contraintes du matériel. Les tâches réalisées les 11 et 12 novembre sont
identifiées séparément des tâches futures.

Le projet n'est pas présenté comme terminé : les fonctions de calcul, les
simulations et la validation expérimentale restent à développer. La prochaine
étape logique sera l'implémentation progressive du prototype Python, en
commençant par la formalisation des types, messages, actions, croyances et
fonctions d'utilité.

\end{document}
